\section{Related literature}\label{sec:literature}

The paper connects several literatures but makes a narrow claim at their intersection.

First, the retail side belongs to the spatial-competition tradition initiated by \citet{hotelling1929}. We use the quadratic endpoint formulation associated with \citet{daspremont1979} in order to keep the conventional differentiation force transparent. The novelty is not a new transport-cost function or an endogenous location result; retailer locations are fixed throughout.

Second, a large literature studies network, clientele, or social externalities in differentiated-price competition. \citet{grilo2001} combine consumption externalities with spatial duopoly and show that conformity can intensify price competition and, when sufficiently strong, generate multiple equilibria. \citet{grivavettas2011} analyze differentiated-product price competition with product-specific network effects and price-sensitive expectations. \citet{tolottiyepez2020} fully characterize Hotelling--Bertrand outcomes with firm-specific network effects, including multiplicity and induced monopoly regimes. More recently, \citet{toshimitsu2026} studies price and quantity competition in a Hotelling market with network connectivity. These papers are the closest reduced-form threat to our pricing mechanism. Our local retail game can indeed be represented by a nonlinear market-share term. We therefore do not claim that market-share feedback, intensified competition, multiplicity, or asymmetric outcomes are new.

Nor is the sign pattern of the price reactions itself new. \citet{tombak2006} defines strategic asymmetry as a situation in which one firm regards the rival's strategic variable as a complement while the other regards its rival's strategic variable as a substitute. In transportation, \citet{mchardy2024} likewise shows that network pricing can contain strategic-complement and strategic-substitute forces simultaneously. Our contribution is therefore not the existence of mixed strategic relations in a transport setting. It is the particular downstream mechanism through which a third-party shared service resource generates the one-sided sign pattern and supports it in a global retail-price equilibrium.

Third, transportation economics already contains rich vertical-market models in which upstream infrastructure or service quality affects downstream competition. \citet{bassozhang2007} analyze congestible facilities with downstream carriers and show that facility charges and capacity/service quality interact with downstream market power. \citet{depalma2018} study transport facilities and downstream carriers in a sequential spatial model with geographical and schedule differentiation. \citet{alvarez2020} analyze an air-transport network in which airport charges affect airline frequencies and fares under asymmetric valuation of frequencies. These papers make clear that an upstream transport layer affecting downstream prices is not new. The distinction here is that the upstream operator does not set a strategic access charge or capacity: conditional on downstream demand it reallocates one fixed shared service resource across rival directions, so a retailer's demand gain changes both its own access and the rival's access.

The shared-resource feature also has an analogue outside transportation. For example, \citet{chenli2013} show that a capacity-allocation mechanism used by a constrained upstream supplier can reshape downstream retailer competition. This literature is an important conceptual threat because it establishes that scarce upstream capacity allocation can be strategically consequential downstream. The present model differs in the object being allocated and in the feedback loop: service allocation changes consumers' generalized access costs, is itself driven by downstream trip demand, and is subject to a minimum-service constraint that can be slack on path while truncating large off-path reallocations.

Fourth, transit economics has long recognized that passenger demand and service frequency interact. The classic analysis of \citet{mohring1972} highlights scale economies associated with service frequency. \citet{baryosef2013} explicitly model vicious and virtuous demand--frequency cycles and multiple equilibria on a bus line. Minimum headway or service-frequency requirements and fixed-fleet constraints are also standard service-design primitives. Operations research provides concrete mechanisms for directional service reallocation: \citet{furth1985} studies deadheading under directional imbalance for a given fleet, while \citet{cortes2011} integrate short-turning and deadheading. More recent demand-responsive service design, such as \citet{liuouyang2021}, makes vehicle allocation and repositioning across zones explicit.

The distinction of the present paper is therefore the complete strategic chain rather than any one ingredient:
\begin{equation*}
\text{retail price}
\rightarrow \text{shopping demand}
\rightarrow \text{third-party fixed-fleet reallocation}
\rightarrow \text{own and rival access}
\rightarrow \text{shopping demand}
\rightarrow \text{retail best responses}.
\end{equation*}
Directional background demand makes this chain asymmetric. The service floor can be slack at the equilibrium while changing large off-path reallocations. In the prior-art audit underlying the paper, no single predecessor was found that combines this demand-responsive third-party shared-resource continuation with a global Hotelling--Bertrand equilibrium exhibiting one-sided strategic substitutability. We consequently position the paper as a transport/spatial-competition contribution, not as a general theorem about network externalities, vertical transport markets, or capacity allocation.
